% Full-report-only exposition of the maintained K-Bound theory.
% This module expands the proofs and scope boundaries without importing any
% historical extension module or introducing a new empirical authority.

\part{A Long-Form Guide to the Population Frontier}
\label{fr:part-theory-exposition}

\section{Keeping the Estimands Apart}
\label{fr:sec-estimands}

The compact presentation develops two related decision problems.  The first is a population
partial-identification problem.  It asks whether the risk difference between a fixed frozen model
and a fixed candidate has one sign in every target law allowed by a declared class.  The second is
an empirical prediction problem.  It asks whether a fitted interval for the measured benefit on a
named evaluation unit lies strictly on one side of zero.  The two problems share an action
vocabulary, but they do not share an estimand or an uncertainty radius.  This chapter slows down
the notation and proofs so that no empirical interval is mistaken for a population identified set.

Throughout the population analysis, the model pair is held fixed.  If an adapter used an unlabeled
batch to construct $f_a$, that construction has already happened.  Population risk is then taken
over a fresh target draw with both $f_0$ and $f_a$ fixed.  The binary result for zero--one loss also
holds the target input marginal, the source information, the score $s$, and every label-free model
output fixed.  Only the target label kernel is allowed to vary inside the declared class.  These
choices are part of the theorem, not implementation details that may be changed silently.

\subsection{A symbol atlas by inferential level}
\label{fr:subsec-symbol-atlas}

Let
\begin{equation}
  d:=\mu_T(D),
  \qquad
  D:=\{x:f_0(x)\ne f_a(x)\}.
  \label{fr:eq-disagreement-mass}
\end{equation}
The maintained population theorem assumes $d>0$.  When $d=0$, the models agree almost surely and
their zero--one risks are equal; the interesting directional problem disappears.  On $D$, let
$s(x)\in[0,1]$ be the fixed candidate-correctness score and
$\eta_a(x)=\Pr_T(f_a(X)=Y\mid X=x)$.  The four central population symbols are
\begin{align}
  M
    &:=\E[s(X)\mid D]-\tfrac12, \\
  \gamma
    &:=\E[\eta_a(X)-s(X)\mid D], \\
  \beta
    &\ge 0 \quad\text{with the declared restriction }|\gamma|\le\beta, \\
  \Delta
    &:=R_T(f_0)-R_T(f_a).
  \label{fr:eq-population-symbol-chain}
\end{align}
Here $M$ is a margin determined by the fixed score and target input law.  The residual $\gamma$ is
a property of the unobserved target label kernel after conditioning on disagreement.  The number
$\beta$ is an external class parameter.  Finally, $\Delta$ is the population risk reduction.  In
particular, $\beta$ is not an estimator, $\gamma$ is not automatically a source-to-target drift
quantity, and $M$ is not the measured accuracy difference on a finite evaluation cell.

The empirical companion uses a separate collection of symbols.  For a predeclared cell
$\mathcal E_j$ and a higher-is-better score $S$,
\begin{equation}
  \Delta_j^{\mathrm{cell}}
  :=S(f_a;\mathcal E_j)-S(f_0;\mathcal E_j).
  \label{fr:eq-cell-benefit}
\end{equation}
A fitted model produces $\widehat\Delta_j$, and a calibration procedure supplies a nonnegative
radius $\varepsilon_j$.  The interval is
$[\widehat\Delta_j-\varepsilon_j,\widehat\Delta_j+\varepsilon_j]$.  These objects describe a
prediction error for the chosen cell-level score.  They do not estimate $(M,\gamma,\beta)$, and
$\varepsilon$ is not a finite-sample version of $\beta$.

\begin{table}[!t]
\centering
\small
\caption{The same model pair appears at two inferential levels.  The population class and the
cell-prediction interval answer different questions.}
\label{fr:tab-symbol-atlas}
\begin{tabular}{@{}p{0.13\textwidth}p{0.20\textwidth}p{0.57\textwidth}@{}}
\toprule
Symbol & Level & Role \\
\midrule
$d$ & population & target probability of the fixed disagreement region $D$ \\
$M$ & population & label-free mean score margin on $D$ \\
$\gamma$ & population & hidden target disagreement-conditional score residual \\
$\beta$ & population & externally declared upper bound on $|\gamma|$ \\
$\Delta$ & population & expected zero--one risk reduction for the fixed model pair \\
$\Delta^{\mathrm{cell}}$ & empirical & observed score difference on one declared evaluation unit \\
$\widehat\Delta$ & empirical & label-free-at-scoring prediction of that cell difference \\
$\varepsilon$ & empirical & calibrated prediction-error radius for the named cell target \\
$B$ & operational & generic scalar target in the coverage-to-action proposition \\
\bottomrule
\end{tabular}
\end{table}

The generic symbol $B$ is useful because the interval implication is agnostic about the scientific
meaning of its scalar target.  Setting $B=\Delta^{\mathrm{cell}}$ gives a cell-level statement.
Setting $B=\Delta$ would give a population statement only if coverage had actually been justified
for $\Delta$.  Substituting one symbol for the other in prose cannot create that justification.

\subsection{Why the conditioning event matters}
\label{fr:subsec-conditioning}

The residual is conditional on $D$.  Even if a score has a useful calibration property under some
source distribution, that fact need not survive either target transfer or conditioning on the
event that the two models disagree.  Disagreement can select a highly atypical region of the input
space.  A score that behaves well on average may have a different residual there.  This is why the
maintained text calls $\gamma$ a calibration residual rather than treating it as a synonym for
distribution shift.

The distinction is visible from the identity
\begin{equation}
  M+\gamma
  =\E[\eta_a(X)\mid D]-\tfrac12.
  \label{fr:eq-residual-identity}
\end{equation}
The score cancels.  Equation~\eqref{fr:eq-residual-identity} is always true under the definitions;
it does not say that either summand is known.  The label-free part $M$ is fixed by the augmented
evidence law used in the theorem.  The hidden part $\gamma$ records exactly how far the average
candidate correctness on $D$ can depart from that score margin.

Treating $M$ as fixed also deserves care.  In the population theorem, $M$ is part of the augmented
observable.  If a practical system estimates $M$ from a finite batch, then estimation error creates
another uncertainty term.  The exact frontier does not make that term vanish.  A valid practical
analysis must account for that estimation error, for example by enlarging the class or proving a
separate sampling statement.  Conditioning is an alternative only under an explicitly justified
conditional observation and model law; conditioning on a model pair adapted on the same batch is
not justified by the fixed-pair theorem.  None of those steps is obtained merely by reusing the
symbol $M$ for a sample average.

\subsection{A disciplined order for reading a claim}
\label{fr:subsec-claim-order}

Every directional statement can be checked in four passes.  First identify the estimand: population
risk, a cell score, or another scalar.  Second identify the information set: label-free target
outputs, labeled historical cells, target labels, or a frozen calibration artifact.  Third identify
the uncertainty object: a declared class radius, a prediction-error radius, or a sampling-error
bound.  Fourth identify the probability law and denominator: all evaluation units, only committed
units, one time point, or an entire sequence.  A statement that omits one of these items is not yet
specific enough to compare with either the population theorem or the empirical proposition.

This order prevents several common substitutions.  A wide $\varepsilon$ does not demonstrate that
$\beta$ is large.  A narrow empirical interval does not show that the population label kernel lies
in a narrow class.  A low number of observed false actions does not establish interval coverage.
And an interval-coverage premise for a random cell does not imply conditional error control among
the cells on which the gate acts.  The remainder of this part derives these boundaries rather than
leaving them as warnings.

\section{Proof Anatomy of the Exact Population Frontier}
\label{fr:sec-frontier-anatomy}

The frontier follows from three ingredients: cancellation outside disagreement, binary
complementarity on disagreement, and an exact description of the residual values allowed by the
declared class.  Each ingredient has a distinct role.  Separating them shows both why the formula is
simple and why its assumptions are restrictive.

\subsection{Step one: cancel the agreement region}
\label{fr:subsec-cancel-agreement}

Under zero--one loss, define
\begin{equation}
  C_a:=\mathbf 1\{f_a(X)=Y\},
  \qquad
  C_0:=\mathbf 1\{f_0(X)=Y\}.
  \label{fr:eq-correctness-indicators}
\end{equation}
Because $\mathbf 1\{f(X)\ne Y\}=1-\mathbf 1\{f(X)=Y\}$, the pointwise loss difference is
\begin{equation}
  \mathbf 1\{f_0(X)\ne Y\}-\mathbf 1\{f_a(X)\ne Y\}
  =C_a-C_0.
  \label{fr:eq-loss-correctness-difference}
\end{equation}
On $D^c$, the two predictions are equal, so $C_a=C_0$.  The integrand is zero there.  Therefore
\begin{align}
  \Delta
  &=\E[C_a-C_0] \\
  &=d\,\E[C_a-C_0\mid D].
  \label{fr:eq-disagreement-reduction-one}
\end{align}
This cancellation uses only equality of the predictions on $D^c$.  It is not a calibration
argument and does not use the class bound.

\subsection{Step two: use binary complementarity}
\label{fr:subsec-binary-complementarity}

On $D$, the two binary predictions differ, so one predicts zero and the other predicts one.  Since
the label is also binary, exactly one prediction is correct.  Thus
\begin{equation}
  C_0=1-C_a \quad\text{on }D.
  \label{fr:eq-binary-complement}
\end{equation}
Substituting into Equation~\eqref{fr:eq-disagreement-reduction-one} gives
\begin{align}
  \Delta
    &=d\,\E[2C_a-1\mid D] \\
    &=2d\left(\Pr_T(f_a(X)=Y\mid D)-\tfrac12\right).
  \label{fr:eq-benefit-from-correctness}
\end{align}
The tower property identifies the conditional correctness probability with
$\E[\eta_a(X)\mid D]$.  Equation~\eqref{fr:eq-residual-identity} then yields
\begin{equation}
  \boxed{\Delta=2d(M+\gamma).}
  \label{fr:eq-benefit-factorization}
\end{equation}
This is the maintained reduction written with $d=\mu_T(D)$.  Because $d>0$, multiplying by $2d$
does not change a sign.  Hence $\sign(\Delta)=\sign(M+\gamma)$ for every target law in the setup,
whether or not that law satisfies a particular residual budget.

Binary complementarity is the step that fails in a direct multiclass transcription.  With three or
more labels, two different predictions can both be wrong.  Cancellation outside $D$ remains valid,
but $C_0=1-C_a$ need not hold on $D$.  The open-extension discussion later in this part begins from
that exact point rather than carrying the binary factor of two into a setting where it is unjustified.

\subsection{Step three: translate the class into a compatible set}
\label{fr:subsec-compatible-set}

Let
\begin{equation}
  q:=\E[\eta_a(X)\mid D]-\tfrac12=M+\gamma.
  \label{fr:eq-normalized-benefit}
\end{equation}
The conditional correctness probability lies in $[0,1]$, so $q\in[-\tfrac12,\tfrac12]$.
The declared restriction $|\gamma|\le\beta$ gives
$q\in[M-\beta,M+\beta]$.  Both restrictions must hold.  The compatible normalized-benefit set is
therefore
\begin{equation}
  \mathcal Q(M,\beta)
  =[-\tfrac12,\tfrac12]\cap[M-\beta,M+\beta]
  =\left[\max\{-\tfrac12,M-\beta\},
          \min\{\tfrac12,M+\beta\}\right].
  \label{fr:eq-compatible-q}
\end{equation}
Scaling by the positive number $2d$ gives the compatible population-benefit set
\begin{equation}
  \mathcal I_\Delta(M,\beta,d)
  =2d\,\mathcal Q(M,\beta).
  \label{fr:eq-compatible-benefit}
\end{equation}
The clipping by $[-\tfrac12,\tfrac12]$ enforces the elementary feasibility of a correctness
probability.  It is not an optional numerical safeguard.  Omitting it would describe residuals
whose implied average correctness is below zero or above one.

Every point in Equation~\eqref{fr:eq-compatible-q} is attainable in the full declared class.  To
see this, choose any $q$ in that interval and set the candidate-correctness kernel on $D$ to the
constant
\begin{equation}
  \eta_q(x):=\tfrac12+q.
  \label{fr:eq-attaining-kernel}
\end{equation}
It is a valid probability because $q\in[-\tfrac12,\tfrac12]$.  Its residual is
$\gamma_q=q-M$, whose absolute value is at most $\beta$ because
$q\in[M-\beta,M+\beta]$.  Holding the input law, models, score, and label kernel on $D^c$ fixed
then produces a target law in the class with benefit $2dq$.  Thus the interval is an identified set,
not merely an outer bound, for the rich class used by the necessity result.

\subsection{Step four: locate the strict-sign regions}
\label{fr:subsec-strict-regions}

A uniformly positive benefit requires the lower endpoint of
$\mathcal Q(M,\beta)$ to be strictly positive.  Because $-\tfrac12$ is not positive, this happens
exactly when $M-\beta>0$, or $M>\beta$.  Similarly, a uniformly negative benefit requires the upper
endpoint to be strictly negative.  Since $\tfrac12$ is not negative, this occurs exactly when
$M+\beta<0$, or $M<-\beta$.  The maximal strict population rule is consequently
\begin{equation}
  M>\beta\Rightarrow\adapt,
  \qquad
  M<-\beta\Rightarrow\freeze,
  \qquad
  |M|\le\beta\Rightarrow\abstain.
  \label{fr:eq-population-rule}
\end{equation}
The inequalities are strict because the action words make strict directional assertions.  A lower
endpoint equal to zero does not justify the claim that every compatible benefit is positive.  An
upper endpoint equal to zero does not justify the claim that every compatible benefit is negative.

For $|M|<\beta$, the interval contains values on both sides of zero.  One may select a small
$r>0$ satisfying $r<\min\{\beta-|M|,\tfrac12\}$ and use the two constant kernels
$\eta_+(x)=\tfrac12+r$ and $\eta_-(x)=\tfrac12-r$.  Their residuals are $r-M$ and $-r-M$; both have
absolute value below $\beta$.  Under the fixed observation experiment in
Section~\ref{fr:subsec-kernel-construction}, they preserve the label-free evidence law and yield
benefits $2dr$ and $-2dr$.  The ambiguity is exact: no additional sample from the same label-free evidence law can
distinguish these two target worlds.

For $|M|=\beta>0$, opposite nonzero signs need not both be available, but zero is.  Taking
$\gamma=-M$ gives $q=0$ and hence $\Delta=0$.  The class also contains a strict-benefit law on the
side indicated by $M$.  Because one compatible law has zero benefit, the strict statement still
fails uniformly.  Finally, when $M=\beta=0$, the class forces $\gamma=0$ and the benefit is
identically zero.  This last case is not uncertainty about the sign: the value is identified, but
neither strict directional assertion is true.

\subsection{Sufficiency, necessity, and class richness}
\label{fr:subsec-richness}

The outside-band direction needs only the bound $|\gamma|\le\beta$.  If $M>\beta$, then
$M+\gamma\ge M-\beta>0$ for every admissible law; the negative case is symmetric.  This is a
sufficiency statement and remains valid for any narrower subclass obeying the same bound.

The inside-band conclusion is different.  To show that no strict direction is uniformly supported,
the class must contain the evidence-matched kernels used above, including the zero-benefit boundary
law.  A narrower scientific model might exclude them.  For example, a verified structural relation
could tie the target label kernel to an observed quantity and leave only one possible sign.  The
maintained theorem therefore states its necessity and maximality for the full rich class.  It does
not say that every imaginable restricted problem has the same identified set.

This distinction is central to using impossibility results responsibly.  An impossibility theorem
does not prohibit action after stronger information is supplied.  It identifies the information
that the present evidence and class fail to provide.  Escaping the result requires changing an
assumption, observing a new channel, or narrowing the target class in a way that can be defended for
the deployment at hand.

\section{Evidence Fibres and the Exact Audit Floor}
\label{fr:sec-fibres}

The frontier assumes a declared residual budget.  A natural next question is whether the same
label-free information used by the gate can discover a smaller valid budget.  Evidence fibres make
the obstacle precise.  They group together target laws that induce exactly the same distribution
for everything the audit may observe.  An audit can transform its evidence, randomize, and compute
arbitrarily complicated statistics, but it cannot tell apart two laws in the same fibre.

\subsection{The observable object and its fibres}
\label{fr:subsec-fibre-definition}

Let $W$ contain the audit's full information set.  It may include the unlabeled target batch,
predictions from both models, logits, features derived from those outputs, the fixed score, and
source or development data whose law is unaffected by the deployment label kernel.  Let an audit
also use an independent random seed $U$ and return the measurable real-valued output
\begin{equation}
  \widehat\beta=g(W,U)\in[0,\infty).
  \label{fr:eq-randomized-audit}
\end{equation}
Including $U$ makes the argument cover randomized procedures.  It gives the audit no additional
information about which compatible target law generated the observable data.

For a declared class $\mathcal C$, two laws are evidence-equivalent when they induce the same law
of $W$.  If $z$ denotes one such observable law, define
\begin{align}
  \mathcal C(z)
    &:=\{P\in\mathcal C:\Law_P(W)=z\}, \\
  \Gamma_z(\mathcal C)
    &:=\sup_{P\in\mathcal C(z)}|\gamma(P)|.
  \label{fr:eq-fibre-radius}
\end{align}
The first object is an evidence fibre.  The second is the largest absolute residual still possible
within it.  Notice that $z$ indexes a probability law, not merely one realized value of a summary
statistic.  Equal realized summaries would be a weaker relation; the exact impossibility uses equal
observable distributions.

The radius is conditional on both the declared class and the chosen evidence.  Enlarging the class
can enlarge a fibre.  Adding a genuinely informative observable can split a fibre into smaller
ones.  Merely computing a new deterministic feature of $W$ cannot do so, because equal laws of $W$
imply equal laws of every measurable function of $W$.

\subsection{How the label-kernel construction preserves evidence}
\label{fr:subsec-kernel-construction}

Fix the target input marginal, source information, predictors, score, and the label kernel outside
$D$.  On $D$, choose any measurable $\eta:D\to[0,1]$ and define the conditional label law by
assigning probability $\eta(x)$ to $Y=f_a(x)$ and probability $1-\eta(x)$ to $Y=f_0(x)$.  This is a
complete binary conditional distribution because the two predictions are distinct on $D$.
Use one fixed marginal-based observation experiment for all target laws, including a fixed joint
law for any source information; finite iid target sampling is one realization.

Changing $\eta$ in this construction changes target labels but changes neither the distribution of
$X$ nor either model's output at a given $X$.  Under this common observation experiment, changing
the target label kernel preserves the law of all recorded label-free evidence.  At the same time,
the residual becomes
\begin{equation}
  \gamma(\eta)
  =\E[\eta(X)\mid D]-\tfrac12-M.
  \label{fr:eq-kernel-residual}
\end{equation}
Thus one evidence law can coexist with a range of target residuals.  The construction isolates the
missing information: correctness on the target disagreement region.

This is stronger than saying that a finite sample is noisy.  The two worlds can induce identical
laws for samples of every size drawn from the permitted label-free channel.  More data from that
same channel can estimate its distribution more accurately, but it cannot reveal which target
label kernel was paired with it.  The obstacle is identification, not a slow convergence rate.

\subsection{Deriving the audit floor line by line}
\label{fr:subsec-audit-proof}

Suppose the audit is required to cover the true absolute residual throughout one nonempty fibre
whose residual set is bounded (as it is for the stated binary score):
\begin{equation}
  \Pr_P\{\widehat\beta\ge|\gamma(P)|\}\ge1-\delta
  \qquad\text{for every }P\in\mathcal C(z),
  \label{fr:eq-audit-validity}
\end{equation}
where $0\le\delta<1$.  Every law in the fibre gives $(W,U)$ the same distribution.  Therefore
$\widehat\beta=g(W,U)$ has one common distribution, which we denote by $\Pr_z$.

If the supremum in $\Gamma_z(\mathcal C)$ is attained at a law $P^\star$, substituting
$P^\star$ into Equation~\eqref{fr:eq-audit-validity} immediately gives
\begin{equation}
  \Pr_z\{\widehat\beta\ge\Gamma_z(\mathcal C)\}\ge1-\delta.
  \label{fr:eq-audit-floor-attained}
\end{equation}
The supremum need not be attained, so the general proof uses approximation.  Select laws $P_n$ in
the fibre such that $a_n:=|\gamma(P_n)|$ increases to $\Gamma_z(\mathcal C)$.  The events
$A_n:=\{\widehat\beta\ge a_n\}$ form a decreasing sequence.  Uniform validity gives
$\Pr_z(A_n)\ge1-\delta$ for every $n$.  By continuity of probability from above,
\begin{align}
  \Pr_z\{\widehat\beta\ge\Gamma_z(\mathcal C)\}
  &=\Pr_z\left(\bigcap_{n=1}^{\infty}A_n\right) \\
  &=\lim_{n\to\infty}\Pr_z(A_n)
  \ge1-\delta.
  \label{fr:eq-audit-floor-limit}
\end{align}
This is the exact fibre-wise floor.  A valid audit must return at least the unresolved fibre radius
with the same high probability at which it promises to cover every compatible residual.  The
constant procedure $\widehat\beta\equiv\Gamma_z(\mathcal C)$ is valid on that fixed fibre, so the
floor is sharp as an oracle benchmark.  Calling it an oracle benchmark matters: the statement does
not provide an algorithm that learns the fibre index or computes its radius from a finite sample.

For the full class $\mathcal C_\beta$ and the maintained feasibility range
$0\le\beta\le\tfrac12$, the constant-kernel construction attains residuals of magnitude $\beta$
while preserving the evidence law.  Hence
\begin{equation}
  \Gamma_z(\mathcal C_\beta)=\beta.
  \label{fr:eq-declared-radius-equality}
\end{equation}
In this setting the label-free audit cannot uniformly certify a radius below the one already
declared.  If the allowed budget is larger than the maintained range, probability feasibility can
clip the attainable residuals, and Equation~\eqref{fr:eq-declared-radius-equality} must not be
repeated without recalculating the fibre.  The general audit-floor statement remains
Equation~\eqref{fr:eq-audit-floor-limit}, with $\Gamma_z(\mathcal C)$ rather than an assumed
numerical substitute.

\subsection{Why labels from an unrelated domain do not shrink the fibre}
\label{fr:subsec-wrong-domain-labels}

Now add a labeled source or calibration sample $L_S$ to the observable object, so
$W=(W_0,L_S)$.  If the declared class places no restriction linking the law of $L_S$ to the target
label kernel on $D$, hold the full joint observation law of $(W_0,L_S)$ fixed while applying the
target kernel construction.  These worlds then have the same recorded evidence law and can still have different values of
$\gamma$.  The labeled sample may estimate source correctness extremely well; it does not identify
target correctness on a region whose label law is allowed to change independently.

The point is not that historical labels are intrinsically useless.  They become informative after
a relation is stated and defended.  Examples include a sampling design that makes them
representative of the intended target units, a transport model that bounds the change in a
conditional label law, or a stability premise that is itself externally validated.  In each case,
the gain comes from the link, not from the word ``labeled.''  Without such a link, adding $L_S$ to
$W$ leaves the decisive target fibre intact.

This observation also clarifies development-based residual surrogates.  A proxy computed from
source-like cells can be useful for engineering exploration.  It does not become a uniform bound
on a deployment residual unless the class restricts how those two quantities differ.  The audit
theorem asks whether the claimed bound covers every target law in the fibre.  Empirical correlation
on development domains is not, by itself, that class restriction.

\subsection{Information refinement is set inclusion}
\label{fr:subsec-information-refinement}

Suppose a new observable $H$ is added.  At a joint observable law $(z,h)$, the refined fibre obeys
\begin{equation}
  \mathcal C(z,h)\subseteq\mathcal C(z),
  \qquad
  \Gamma_{z,h}(\mathcal C)\le\Gamma_z(\mathcal C).
  \label{fr:eq-fibre-refinement}
\end{equation}
The inequality is simply monotonicity of a supremum under set inclusion.  It gives a clean language
for escape routes.  New information helps only to the extent that it excludes target laws with
large residuals.  If $H$ carries representative target labels or a validated structural
measurement, the inclusion may be strict.

If $H$ is a deterministic function of $W$, then for every compatible joint observable law
$(z,h)$, the refined fibre equals the original fibre:
\begin{equation}
  \mathcal C(z,h)=\mathcal C(z).
  \label{fr:eq-deterministic-refinement}
\end{equation}
Indeed, the law of $(W,H)$ is already determined by the law $z$ of $W$.  A deterministic recoding
can make the evidence easier to use, but it cannot exclude a target law that was compatible with
the original evidence distribution.

Equation~\eqref{fr:eq-fibre-refinement} is deliberately one-sided for genuinely new evidence.  It does not say that any
particular label-shift statistic, rank condition, or symmetry assumption produces a symmetric
interval or an exact if-and-only-if frontier.  Those conclusions require a separate model-specific
proof.  Likewise, a closure property under swapping labels does not automatically make arbitrary
regional benefit signs feasible.  The maintained population theorem obtains exactness from its
explicit binary kernel construction, not from an abstract slogan about symmetry.

\section{Strict Actions, Closed Bands, and Coverage}
\label{fr:sec-action-semantics}

The three action names carry logical content.  They are not merely labels for which model is served.
An adapt decision asserts that the candidate has strictly positive benefit.  A \freeze decision
asserts that the candidate has strictly negative benefit.  An abstain decision makes neither
assertion.  In the implementation described in the main text, both \freeze and abstain retain
$f_0$, but their audit meanings remain different: one is a supported rejection and the other is a
record of unresolved evidence.

\subsection{Why the boundary is closed}
\label{fr:subsec-closed-band}

The population compatible set supports adapt only if every member is greater than zero, and it
supports \freeze only if every member is less than zero.  The word ``every'' explains the closed
abstention band.  At $M=\beta>0$, the smallest compatible normalized benefit is zero.  At
$M=-\beta<0$, the largest is zero.  At $M=\beta=0$, zero is the only compatible benefit.  In all
three cases a strict statement is false for at least one admissible law.

This convention is intentionally stronger than choosing a model with zero regret.  When
$\Delta=0$, either model has the same population risk, so serving either one incurs no loss relative
to the pairwise oracle.  Yet neither ``the candidate is better'' nor ``the frozen model is better''
is true.  The frontier concerns warranted directional assertions, not only realized task loss.

The same logic applies to an empirical interval.  Let
\begin{equation}
  L:=\widehat B-\varepsilon,
  \qquad
  U:=\widehat B+\varepsilon.
  \label{fr:eq-operational-endpoints}
\end{equation}
The operational rule uses $L>0$ for adapt and $U<0$ for \freeze.  If $L=0$ or $U=0$, the interval
touches the boundary and the rule abstains.  Replacing either strict comparison by a weak one would
change the claim: it would permit a directional action even though the interval contains a value
at which the corresponding strict assertion fails.

Unavailable evidence is another boundary case.  If required features, a candidate, or a frozen
calibration artifact are unavailable, there is no interval on which to base a negative claim.
Retaining $f_0$ may be the appropriate service behavior, but its audit record is abstain, together
with the reason.  Reclassifying missing assessment as \freeze would make the apparent negative
action count larger without supplying negative evidence.

\subsection{What exact matched evidence forces a randomized rule to do}
\label{fr:subsec-randomized-actions}

Consider two compatible population laws with exactly the same evidence distribution and opposite
nonzero benefit signs.  A possibly randomized label-free rule has the same action probabilities in
both worlds.  If its marginal false-adapt probability may not exceed $\alpha$ in any world, then
its probability of adapt is at most $\alpha$ in the negative-benefit world.  Because the action
distribution is shared, the same upper bound applies in the positive-benefit world.  Symmetrically,
control of false \freeze limits its \freeze probability to at most $\alpha$.  The remaining mass
must satisfy
\begin{equation}
  \Pr\{g=\abstain\}\ge 1-2\alpha.
  \label{fr:eq-matched-abstention}
\end{equation}
Randomization therefore cannot manufacture information.  It can distribute a permitted error
budget across unsupported actions, but a rule that controls both directions must leave most mass on
abstention when the evidence law is shared by opposite worlds.  This is a probabilistic error-control
statement.  It is distinct from the deterministic pointwise-maximal rule, which abstains everywhere
that no strict action is uniformly sound.

\subsection{Coverage-to-action as a set containment}
\label{fr:subsec-coverage-containment}

The operational proposition has an even shorter proof because it begins from a coverage premise.
Let
\begin{equation}
  E:=\{|\widehat B-B|\le\varepsilon\}.
  \label{fr:eq-coverage-event}
\end{equation}
On $E$, an adapt decision implies
$B\ge\widehat B-\varepsilon>0$.  Therefore
\begin{equation}
  \{g=\adapt,\ B\le0\}\subseteq E^c.
  \label{fr:eq-false-adapt-containment}
\end{equation}
Likewise, on $E$, a \freeze decision implies
$B\le\widehat B+\varepsilon<0$, and hence
\begin{equation}
  \{g=\freeze,\ B\ge0\}\subseteq E^c.
  \label{fr:eq-false-freeze-containment}
\end{equation}
If $\Pr(E)\ge1-\alpha$, each marginal false-direction probability is at most $\alpha$.  No
distributional approximation appears in these two inclusions.  All statistical difficulty has
been placed in the premise that the interval covers the predeclared target under the intended joint
law.

The proposition does not prove its own premise.  An order statistic, a regression algorithm, or the
name of a calibration method does not by itself establish coverage for a changed candidate,
environment, or evaluation unit.  One must specify the unit distribution, all fitting and candidate
construction randomness, and the assumptions under which calibration units and a future unit are
related.  The action implication is exact conditional on that work being valid; it cannot replace
that work.

\subsection{Marginal error, conditional error, and exposure}
\label{fr:subsec-error-denominators}

Equation~\eqref{fr:eq-false-adapt-containment} bounds
$\Pr(g=\adapt,B\le0)$, whose denominator is the entire probability space.  It does not directly
bound
\begin{equation}
  \Pr(B\le0\mid g=\adapt),
  \label{fr:eq-conditional-false-adapt}
\end{equation}
whose denominator is only the adapt actions.  When $\Pr(g=\adapt)>0$, the marginal statement gives
the elementary ratio bound
\begin{equation}
  \Pr(B\le0\mid g=\adapt)
  \le \min\left\{1,\frac{\alpha}{\Pr(g=\adapt)}\right\}.
  \label{fr:eq-exposure-ratio}
\end{equation}
This can be uninformative when action exposure is small.  Obtaining a useful fixed level conditional
on selection requires a different guarantee.  If the rule never adapts, its marginal false-adapt
probability is zero mechanically, while the conditional rate is undefined.  Thus every safety
summary must report exposure alongside error.

The same denominator discipline applies to finite recorded panels.  A count divided by all cells
answers a marginal descriptive question under equal cell weighting.  A count divided by actions of
one type answers a conditional descriptive question on the exposed subset.  Neither ratio is an
independent estimate for a new environment unless the sampling and dependence structure support
that interpretation.

\subsection{Choosing the target \texorpdfstring{$B$}{B}}
\label{fr:subsec-choosing-B}

For the empirical gate, $B=\Delta^{\mathrm{cell}}$ yields a statement about the named cell-level
score.  This target may be an accuracy difference, a macro-averaged score difference, or another
predeclared scalar.  The population theorem instead concerns $B=\Delta$, the expected zero--one
risk reduction for a fresh target draw and a fixed binary model pair.  Coverage of the first does
not entail coverage of the second.

A bridge can be built only by adding an explicit relation.  For example, suppose one event gives
$|\widehat\Delta-\Delta^{\mathrm{cell}}|\le\varepsilon$ and another gives
$|\Delta^{\mathrm{cell}}-\Delta|\le b$.  On their intersection, the triangle inequality gives
\begin{equation}
  |\widehat\Delta-\Delta|\le\varepsilon+b.
  \label{fr:eq-cell-population-bridge}
\end{equation}
A probability statement then needs valid coverage levels for both events and an appropriate union
bound or joint argument.  Without the second event, the extra radius $b$, and their assumptions,
the empirical interval remains a cell-benefit interval.

Selection and time introduce further distinctions.  Marginal coverage for each of several
candidates does not automatically cover the candidate selected after inspecting their intervals.
Pointwise coverage at each time does not control the probability of any error over an indefinite
sequence.  Reusing outcomes to tune the predictor or radius changes the probability space.  These
are not defects in the elementary containment proof; they are different statistical targets that
need stronger premises.

\section{What Must Change to Escape the Impossibility}
\label{fr:sec-escape}

An escape route must exclude some of the evidence-matched label kernels.  That can happen through a
new observation or through a narrower class, but the new restriction must be stated at the same
target level as the desired conclusion.  A useful way to audit a proposed route is to ask three
questions: What target laws does it remove from the fibre?  Why is that removal credible in the
deployment domain?  What finite-sample uncertainty remains after the structural restriction is
accepted?

\subsection{Representative target labels}
\label{fr:subsec-target-labels}

Labels sampled from the intended target-unit distribution directly reveal target correctness.  In
the binary setup, labeled observations on $D$ can estimate
$\E[\eta_a(X)\mid D]$ and hence $M+\gamma$.  Their usefulness depends on disagreement exposure: if
$D$ is rare, a general target sample can contain few informative labels.  A sampling design that
targets disagreements may be more efficient, but then its weighting and selection mechanism must
be accounted for.  In either design, labels replace an identification gap with a sampling problem;
they do not make uncertainty disappear.

Historical deployment labels can play the same role when the past and future units are connected by
a defensible sampling or stability model.  The resulting statement is episode- or environment-
specific.  If conditions can change after calibration, that change requires its own bound or
monitoring assumption.  The evidence fibre shrinks because the class couples past and future target
kernels, not because past labels possess universal transfer value.

\subsection{Structural coupling across domains}
\label{fr:subsec-domain-coupling}

A model may restrict how the target label kernel differs from a source or calibration kernel.  A
uniform bound on an appropriate conditional discrepancy, a verified invariance, or a mechanistic
constraint can imply a residual budget.  Such a route is scientifically meaningful when the
constraint has a domain interpretation and can be stress-tested independently.  Simply naming the
observed change ``covariate shift'' does not establish the needed equality of conditional label
laws.

Label-shift models offer another possible restriction: fixed class-conditional feature laws plus a
changing class prior can make some target quantities recoverable through an evidence operator.  At
the abstract level, adding informative restrictions cannot enlarge a fibre, as
Equation~\eqref{fr:eq-fibre-refinement} shows.  Stronger conclusions---including a symmetric
identified interval, an exact sign frontier, or an if-and-only-if condition---depend on the
particular operator, feasibility constraints, and invertibility assumptions.  They remain open here
and are not consequences of monotonicity alone.

\subsection{Additional views and proxy assumptions}
\label{fr:subsec-additional-views}

Multiple predictions, augmentations, sensors, or independently trained models can supply a richer
observable $H$.  They help only under a relation between their agreement pattern and correctness.
Conditional-independence or error-structure assumptions can sometimes provide such a relation, but
agreement alone is not truth.  A rank pattern may be a useful diagnostic without being a two-sided
certificate of latent error.  In particular, a small observable residual does not upper-bound an
unobservable model deviation when the proved inequality goes in the opposite direction.

A domain expert may instead declare a physically meaningful upper bound on the residual.  This is
not a statistical estimate, but it can be the right object when the mechanism supplies knowledge
that the data channel cannot.  The decision is then conditional on that declared class.  A good
report should vary the bound in a sensitivity analysis and explain what evidence would falsify the
assumption, while avoiding the claim that the unlabeled batch verified it.

\subsection{An assumption ledger for any proposed escape}
\label{fr:subsec-assumption-ledger}

Before promoting an escape route to a decision rule, its ledger should name: the target population;
the fixed or adaptively constructed model pair; the exact information available at action time; the
new structural restriction; the parameter or residual it bounds; the unit over which any coverage
statement is marginal; and the behavior when a premise cannot be checked.  It should also separate
identification from estimation.  The first asks what would be known with the evidence distribution
known exactly.  The second asks how accurately a finite sample recovers the identified object.  A
rate calculation cannot repair non-identification, while an identification argument does not supply
a finite-sample rate.

\section{Open Extensions Beyond the Maintained Binary Pair}
\label{fr:sec-open-extensions}

The following directions are research programs, not additional results claimed by this report.
Each starts from an identity that remains valid and then states the missing ingredient needed for a
frontier or certificate.  This format prevents an attractive analogy from being presented as a
proved extension.

\subsection{Multiclass prediction}
\label{fr:subsec-multiclass-open}

For multiclass zero--one loss, cancellation outside $D$ still gives
\begin{equation}
  \Delta
  =d\left(p_a-p_0\right),
  \quad
  p_a:=\Pr_T(f_a(X)=Y\mid D),
  \quad
  p_0:=\Pr_T(f_0(X)=Y\mid D).
  \label{fr:eq-multiclass-reduction}
\end{equation}
The difference from the binary proof is exact: two distinct multiclass predictions may both be
wrong, so $p_0$ need not equal $1-p_a$.  A useful generalization would model the conditional
correctness advantage
\begin{equation}
  A(x):=\Pr_T(f_a(X)=Y\mid X=x)
        -\Pr_T(f_0(X)=Y\mid X=x)
  \label{fr:eq-multiclass-advantage}
\end{equation}
and compare it with a label-free proxy $u(x)$.  Defining a proxy mean and a residual yields an
algebraic decomposition of $\Delta$, and an externally justified residual bound gives an immediate
outside-band sufficiency statement.  An exact necessity result would require characterizing which
advantage functions can arise from one multiclass label kernel while all model outputs remain
fixed.  The feasibility geometry is not the binary interval, so the maintained binary theorem does
not settle it.

\subsection{Regression and general losses}
\label{fr:subsec-regression-open}

For a general loss, define the random loss advantage
\begin{equation}
  G(X,Y):=\ell(f_0(X),Y)-\ell(f_a(X),Y),
  \qquad
  \Delta=\E[G(X,Y)].
  \label{fr:eq-general-loss-advantage}
\end{equation}
If a label-free proxy $t(X)$ is available, the identity
\begin{equation}
  \Delta=\E[t(X)]+\E[G(X,Y)-t(X)]
  \label{fr:eq-general-loss-decomposition}
\end{equation}
suggests a margin-plus-residual analysis.  A declared absolute bound on the second term is sufficient
for a strict sign outside the corresponding band.  Exactness is harder.  It depends on the loss,
the allowed response distributions, integrability, output bounds, and which label kernels preserve
the evidence.  Unbounded losses also require tail conditions before a finite residual budget is
meaningful.  Previously suggested regression brackets and constants are not imported here; a valid
extension needs a fresh statement and proof under explicit assumptions.

\subsection{More than one candidate}
\label{fr:subsec-multiple-candidates-open}

With candidates $f_1,\ldots,f_K$, define each benefit relative to the same frozen reference,
$B_k=R_T(f_0)-R_T(f_k)$ or its cell-level analogue.  Marginally valid intervals for the $K$
predeclared targets do not automatically remain valid after selecting the candidate with the most
favorable predicted benefit.  A sound design could reserve independent information for selection
and evaluation, construct a simultaneous event for all candidates, or use a procedure proved valid
conditional on the selection rule.  Each choice changes exposure and power.

The population identification problem also changes.  Pairwise model agreements constrain a shared
label kernel, and the feasible vector $(B_1,\ldots,B_K)$ need not be a product of pairwise intervals.
Low-rank agreement patterns can suggest structure, but a one-sided fit diagnostic is not an upper
certificate for violation of that structure.  Nor does pairwise agreement by itself establish
conditional independence of correctness indicators.  A complete theory would characterize the
joint feasible benefit set, state the sign anchor that orients it, and prove how an observable test
controls departures from the assumed model.  Those steps remain open.

\subsection{Sequential and adaptive deployment}
\label{fr:subsec-sequential-open}

In a sequence, let $\mathcal F_t$ contain all information available before action $g_t$.  The
candidate at time $t$ may depend on earlier inputs and actions, and an accepted update may become the
next baseline.  The target $B_t$, its predictor, and its calibration distribution can therefore all
change with the history.  A marginal guarantee at each fixed time does not imply
\begin{equation}
  \Pr\{\text{some false directional action occurs over time}\}\le\alpha.
  \label{fr:eq-sequential-family-event}
\end{equation}
Controlling this event requires a sequential design, such as a valid allocation of error across
times, a time-uniform confidence sequence, or another supermartingale-based construction, together
with assumptions robust to the feedback created by adaptation.  Recalibrating on outcomes selected
because an earlier gate was uncertain can also induce selection bias.

The population fibre becomes history-dependent as well.  New labels can shrink it, while model
updates and distribution change can redefine the disagreement region and residual.  A sequential
extension should specify when the declared class is refreshed, what remains fixed between actions,
and whether safety is required per action, over a finite horizon, or uniformly in time.  The present
theory and ordinary cell-level interval do not provide that guarantee.

\subsection{A research checklist for extensions}
\label{fr:subsec-extension-checklist}

An extension is ready to join the maintained theorem stack only after six items are explicit:
\begin{enumerate}
  \item the estimand and action semantics;
  \item the observable evidence and its probability law;
  \item the admissible target class and feasibility constraints;
  \item an exact or conservative compatible-value set with a proved construction or bound;
  \item the statistical premise needed to estimate or cover the intended target; and
  \item the denominator, selection rule, and time horizon of every error statement.
\end{enumerate}
Numerical checks may reveal counterexamples and test algebra, but they do not substitute for these
items.  Conversely, a population proof does not certify a software pipeline or empirical transfer.
The value of the K-Bound decomposition is that it makes this division of labor visible: structural
assumptions determine what can be known, statistical calibration governs finite-sample uncertainty,
and action semantics determine which claims are made when either layer is inconclusive.
